% !TeX root = main.tex
\section{Introduction}
Flapping-wing micro air vehicles (FWMAVs) operate in a regime where aerodynamic load generation is tightly coupled to kinematics, body attitude, and structural compliance.
In practice, control authority is often realized by adjusting flapping frequency, stroke amplitude, and/or phase relationships between multiple joints or wings.
However, these inputs are not always independently adjustable due to actuator limits, joint constraints, and the need to maintain stable periodic wing motion.

Duty-cycle modulation changes the fraction of the cycle spent on downstroke versus upstroke while keeping the stroke amplitude fixed.
Mechanically, this can be implemented by time-reparameterizing a periodic stroke trajectory.
Because aerodynamic loads depend nonlinearly on instantaneous velocity and on phase-dependent effective angle of attack, a time-reparameterization can significantly alter cycle-averaged thrust, lift, and power.
This makes duty-cycle modulation an attractive candidate for a low-dimensional control input, especially for forward-speed regulation where the thrust component is directly relevant.

This paper evaluates duty-cycle modulation in a controlled simulation environment and connects the observed trends to control-relevant tradeoffs near a trim-like operating condition.
We focus on the downstroke ratio parameter $\delta \in (0,1)$, quantify the mapping $\delta \mapsto (\bar{F}_x, \bar{F}_z, \bar{P}_{in})$, and summarize non-dominated operating points via a Pareto analysis.
We also discuss implementation on a physical platform and present a short closed-loop forward-speed regulation demonstration using duty-cycle modulation, together with practical insights about actuator bandwidth, compliance, and model mismatch.

\subsection{Contributions}
The contributions of this work are:
\begin{itemize}
\item A reproducible implementation of split-cycle duty-cycle modulation (SCCPFM) in an IsaacLab flapping-wing wind-tunnel setup.
\item A trim-based sweep of downstroke ratio quantifying cycle-averaged thrust, lift, and power tradeoffs, including Pareto-front extraction under lift feasibility constraints.
\item A practical discussion and short closed-loop validation of duty-cycle modulation as a forward-speed control input on a physical flapping platform.
\end{itemize}
